\section{Limitations, Ethics, and Reuse}
\label{sec:limitations}

\lafc's labels and evaluation have several important constraints, and the
practical guidance in Section~\ref{sec:practical-implications} should be
read together with them rather than instead of them.

First, the main label is a finite-horizon, continuation-policy-specific
counterfactual miss cost, not an offline-optimal label for the full future
control problem. Good performance on this label does not automatically
imply online deployment gains, which is exactly why
Sections~\ref{sec:closed-loop} and~\ref{sec:linkage} evaluate closed-loop
behavior directly rather than relying on the offline label alone.

Second, horizon-dependent results are scoped to the horizons actually
evaluated. Most manuscript analyses use the canonical released
$H\in\{4,8,16\}$ decision population; the matched long-horizon audit extends
only the discriminativeness characterization through
$H\in\{32,64,128\}$ on the same physical decision population.
It shows that longer horizons increase informativeness but do not remove
degeneracy: substantial all-tied mass persists at $H=128$, and wiki2018
remains fully tied. No result here covers $H>128$.

Third, the canonical labels use LRU continuation after a forced eviction.
The sampled MRU/random study and full-population MRU census in
Section~\ref{sec:continuation} support robustness only within their stated
scopes; random continuation is not validated at population scale, and no
experiment covers every deployment policy.

Fourth, the target is substantially tied (Section~\ref{sec:characterization}:
all-tied fraction 0.68, unique-winner fraction exactly 0.0). Support for
the benchmark's usefulness is conditional on decision informativeness, not
universal, and no result in this paper claims it generalizes beyond the
audited population. wiki2018 remains a fully degenerate negative control
at every capacity and horizon tested; Section~\ref{sec:mechanisms} proves
this follows deductively from zero reuse events, not as a statistical
property of the trace.

Fifth, the evaluated dataset uses unit-object, count-based capacities (32,
64, 128, or 256 objects) with unconditional admission of the requested
object. Variable object sizes, byte-weighted miss cost, latency-aware cost,
and dirty-write cost are not modeled by this label family, and admission
is never withheld. Claims in this paper should be read as scoped to this
abstraction and not generalized to byte-addressable or admission-gated
cache models without further evidence.

Sixth, closed-loop evidence carries two family-specific caveats stated in
full in Section~\ref{sec:experimental-methodology}: MetaCDN's evidence is
validation-window, not test-window; MetaKV's scored window begins at an
empty cache. Section~\ref{sec:mechanisms} discusses, without resolving,
whether the cold start plausibly contributes to MetaKV being the one
cell where LRU is not the best policy.

Seventh, a known data-generation issue remains open: in the audited
canonical data, the predictor-victim flag is identical to the LRU-victim
flag throughout, and related predictor, bucket, and confidence fields
include constant or placeholder values. This paper does not make any claim
that depends on those predictor fields being informative.

Eighth, benchmark reuse must preserve provenance distinctions. The release
exposes generated supervision and derived views, but it does not erase the
attribution, citation, or licensing requirements of the upstream raw traces
from which candidate rows were derived; those raw traces remain external
artifacts with their own governance constraints. The open subset is also
intentionally narrower than the full generated corpus for redistribution
reasons, and not every selected family appears in every split of the
evaluated release package. The exact generator commit for the full
historical release corpus is not completely reconstructed; this is
orthogonal to the audited evidence reported in this paper but is disclosed
here for reproducibility completeness.

Finally, this paper reports closed-loop results for LRU, MRU, random, and
SIEVE (Section~\ref{sec:closed-loop}), plus ARC, LIRS, and S3-FIFO under
the identical protocol (Section~\ref{sec:expanded-comparators}); it does
not report a closed-loop evaluation of any learned policy trained on these
labels, and the seven-policy set is still representative rather than
exhaustive -- CLOCK-Pro, GreedyDual-Size, and TinyLFU-family admission
variants beyond W-TinyLFU (excluded after a validation failure at capacity
32; see Section~\ref{sec:expanded-comparators}) remain untested here. A frozen pilot
(not the production Tier-1 evaluation reported here) previously found a
simple learned policy competitive with random and SIEVE but not with LRU
on one family, with the other pilot family blocked on the absence of an
independent test window; neither result is repeated or extended here, and
learned-policy closed-loop competitiveness remains future work rather than
a claim of this paper.
